\clearpage\pagestyle{otMain}\setlength{\otBlocEspacement}{2.5mm}\fontsize{11}{15}\selectfont\corrigesChapitre{6}{Équations et résolution de problèmes}
\begin{blocLivre}{cor-6-1}\begin{corrige}{1}{Tester une valeur}\demarche{Substituer sans transformer l’équation.}
Pour 4, $3\times4-2=10$ : $\res4$ est solution. Pour 5, $3\times5-2=13\ne10$ : 5 n’est pas solution. Tester consiste à calculer les deux membres.
\end{corrige}\end{blocLivre}
\begin{blocLivre}{cor-6-2}\begin{corrige}{2}{Une addition}\demarche{Effectuer la même soustraction dans les deux membres.}
On soustrait 8 aux deux membres : $x=-3-8=\res{-11}$. Vérification : $-11+8=-3$.
\end{corrige}\end{blocLivre}
\begin{blocLivre}{cor-6-3}\begin{corrige}{3}{Un coefficient négatif}\demarche{Diviser les deux membres par le coefficient de l’inconnue.}
On divise par $-4$, nombre non nul : $x=18/(-4)=\res{-9/2}$, soit $-4{,}5$. Le produit $(-4)\times(-4{,}5)=18$ confirme.
\end{corrige}\end{blocLivre}
\begin{blocLivre}{cor-6-4}\begin{corrige}{4}{Deux étapes}\demarche{Enlever la constante avant le coefficient.}
On ajoute 7 : $5x=30$. On divise par 5 : $\res{x=6}$. Contrôle : $5\times6-7=23$.
\end{corrige}\end{blocLivre}
\begin{blocLivre}{cor-6-5}\begin{corrige}{5}{Traduire une phrase}\demarche{Distinguer le triple du nombre et l’ajout final.}
Si le nombre est $x$, la phrase donne $3x+4=25$. Donc $3x=21$ puis $\res{x=7}$. Contrôle : $3\times7+4=25$.
\end{corrige}\end{blocLivre}
\begin{blocLivre}{cor-6-6}\begin{corrige}{6}{L’inconnue des deux côtés}\demarche{Réunir les termes en $x$ dans un membre.}
On soustrait $3x$ : $4x-5=11$. On ajoute 5 : $4x=16$, puis $\res{x=4}$. Les deux membres valent 23 pour $x=4$.
\end{corrige}\end{blocLivre}
\begin{blocLivre}{cor-6-7}\begin{corrige}{7}{Avec une parenthèse}\demarche{Développer avant de regrouper.}
On développe : $3x-6=x+8$. On soustrait $x$ et ajoute 6 : $2x=14$, donc $\res{x=7}$. Contrôle : $3(7-2)=15=7+8$.
\end{corrige}\end{blocLivre}
\begin{blocLivre}{cor-6-8}\begin{corrige}{8}{Avec des fractions}\demarche{Simplifier les constantes avant de multiplier.}
On soustrait $1/2$ : $x/3=4/2=2$. On multiplie par 3 : $\res{x=6}$. Contrôle : $6/3+1/2=5/2$.
\end{corrige}\end{blocLivre}
\begin{blocLivre}{cor-6-9}\begin{corrige}{9}{Corriger une méthode}\demarche{Appliquer une opération au membre entier.}
En divisant par 2, il faut diviser tout le membre : $x+3=7$, donc $\res{x=4}$. On peut aussi soustraire 6 d’abord : $2x=8$. Le contrôle donne $2\times4+6=14$.
\end{corrige}\end{blocLivre}
\begin{blocLivre}{cor-6-10}\begin{corrige}{10}{Aucune ou toutes}\demarche{Observer ce qu’il reste après réduction.}
La première donne $2=-1$, impossible : aucune solution. La seconde se développe en $3x+6=3x+6$, vraie pour tout nombre réel : toutes les valeurs réelles conviennent.
\end{corrige}\end{blocLivre}
\begin{blocLivre}{cor-6-11}\begin{corrige}{11}{Un rectangle}\demarche{Traduire le périmètre avec la même inconnue.}
$2x+2(x+5)=38$, soit $4x+10=38$. Donc $4x=28$ et $\res{x=7}$. La largeur est 7 cm et la longueur $\res{12}$ cm. Contrôle : $2(7+12)=38$ cm, avec deux dimensions positives.
\label{fig-cor-6-11}\figureGeo[corrige]{eq-rectangle}
\end{corrige}\end{blocLivre}
\begin{blocLivre}{cor-6-12}\begin{corrige}{12}{Deux tarifs}\demarche{Égaler les deux expressions de coût.}
Pour $x$ entrées : $12+3x=5x$. Donc $12=2x$, soit $\res{x=6}$. Les deux dépenses valent alors $\res{30}$ euros. Le nombre d’entrées obtenu est bien entier et positif.
\end{corrige}\end{blocLivre}
\begin{blocLivre}{cor-6-13}\begin{corrige}{13}{Des entiers consécutifs}\demarche{Exprimer tous les nombres avec une inconnue.}
On note $n$ le premier ; les deux suivants sont $n+1$ et $n+2$. $3n+3=81$ donne $3n=78$ puis $n=26$. Les nombres sont $\res{26,27,28}$ et leur somme est bien 81.
\end{corrige}\end{blocLivre}
\begin{blocLivre}{cor-6-14}\begin{corrige}{14}{Un âge futur}\demarche{Faire augmenter les deux âges du même nombre d’années.}
Dans $x$ années, les âges seront $12+x$ et $40+x$. On écrit $40+x=2(12+x)$. Donc $40+x=24+2x$ et $\res{x=16}$. Les âges seront 28 et 56 ans : le second est bien le double du premier.
\end{corrige}\end{blocLivre}
\begin{blocLivre}{cor-6-15}\begin{corrige}{15}{Une longueur admissible}\demarche{Résoudre puis contrôler les contraintes de longueur.}
$3(x+2)=4x$ donne $3x+6=4x$, donc $\res{x=6}$. Le carré a pour côté 6 cm et le triangle 8 cm : les deux longueurs sont positives. Les périmètres valent chacun $\res{24}$ cm.
\label{fig-cor-6-15}\figureGeo[corrige]{eq-deux-figures}
\end{corrige}\end{blocLivre}
